\section{不定式与L'Hôpital法则}
\label{sec:ru-lhopital}

两个都趋于零的函数相除，比值的极限不能由各自的极限算出，这类形式称为不定式。它问的其实是两个无穷小谁更快趋于零。
函数在趋近点附近可微时，两个函数的增量都可以用导数描述，于是可以借导数比较它们。

\begin{theorem}[L'Hôpital法则]
\label{thm:ru-lhopital}
设$f,g$在$(a,a+r)$内可导，$g'\neq0$，并且$x\to a^+$时$f(x)\to0$、$g(x)\to0$。若$f'(x)/g'(x)\to L$，其中$L$为实数或$\pm\infty$，则$f(x)/g(x)\to L$。左侧与双侧极限有相同的结论。
\end{theorem}
\begin{proof}
补充定义$f(a)=g(a)=0$，则$f,g$在每个$[a,x]$（$x<a+r$）上连续、在内部可导。由Cauchy中值定理的比值形式，存在$\xi_x\in(a,x)$使
\[
 \frac{f(x)}{g(x)}=\frac{f(x)-f(a)}{g(x)-g(a)}=\frac{f'(\xi_x)}{g'(\xi_x)}.
\]
$x\to a^+$时$\xi_x\to a^+$，右侧趋于$L$。比值形式所需的$g(x)\neq g(a)$由$g'\neq0$保证。
\end{proof}
定理中的条件都是充分条件：在趋近点的一侧可导、分母的导数非零、导数比有极限。例如
\[
 \lim_{x\to0}\frac{e^x-1-x}{x^2}=\lim_{x\to0}\frac{e^x-1}{2x}=\frac12,
\]
第二个极限仍是$0/0$型，它由基本极限$(e^x-1)/x\to1$直接得到。法则对$x\to\pm\infty$以及分子分母都趋于无穷的$\infty/\infty$型也成立，证明思路相同\scite{math-lebl-analysis1}。

其他不定式要先转换成$0/0$或$\infty/\infty$，再检查条件。$x\to0^+$时$x\log x$是$0\cdot\infty$型，写成$\log x/(1/x)$后为$\infty/\infty$型，导数比$(1/x)/(-1/x^2)=-x\to0$，所以$x\log x\to0$。
$1^\infty$与$0^0$型先取对数：$(1+1/n)^n=e^{n\log(1+1/n)}$，而$n\log(1+1/n)=\log(1+t)/t$（$t=1/n$）趋于$\log$在$1$处的导数$1$，由指数的连续性，$(1+1/n)^n\to e$。

法则不是万能的，有三点要注意。第一，导数比的极限不存在时，原极限仍可能存在：$x\to+\infty$时$(x+\sin x)/x\to1$，导数比$1+\cos x$却在$0$与$2$之间振荡。
第二，夹逼和等价无穷小比较仍然是独立的工具，很多时候更直接。第三，不能用L'Hôpital法则回头证明$\sin h/h\to1$或$(e^h-1)/h\to1$：
对前者使用法则需要$(\sin)'(0)=1$，而这个导数正是从该极限算出来的，论证会循环。基本极限只能像第~\ref{sec:ru-elementary-limits}节那样用几何或级数直接估计。
